\section{C++ Library}
\label{sec:CppLibrary}

The C++ implementation lives under \filepath{graft-cpp}. It is built separately
with CMake so C++ deployments do not depend on the Maven reactor. It shares the
protobuf schema and envelope protocol with Java.

\subsection{Build model}

The C++ build creates:

\begin{itemize}
\item \texttt{graft\_proto}: generated protobuf types,
\item \texttt{graft\_transport}: reusable library target for core/runtime/storage/transport,
\item \texttt{graft\_smoke}: command-line probe and runnable modes,
\item \texttt{graft\_unit\_tests}: Catch2 tests when Catch2 v3 is available.
\end{itemize}

Basic workflow:

\begin{lstlisting}[language=bash,caption={C++ build workflow}]
cmake -S graft-cpp -B graft-cpp/build
cmake --build graft-cpp/build
ctest --test-dir graft-cpp/build --output-on-failure
\end{lstlisting}

The code requires C++20. Older Boost.Asio installations may still reference
\texttt{std::result\_of}; the CMake target defines
\texttt{BOOST\_ASIO\_DISABLE\_STD\_RESULT\_OF} to remain compatible with such
headers in C++20 mode.

\subsection{Application boundary}

The C++ application boundary mirrors the Java state-machine boundary:

\begin{itemize}
\item commands are opaque bytes or strings at the Raft layer,
\item application state machines apply committed commands,
\item snapshots are application bytes wrapped by Raft metadata,
\item queries are served only after the runtime establishes read safety.
\end{itemize}

The current C++ example state machine is key-value oriented. Reference-data mode
behavior exists at the policy boundary, but the Java implementation remains the
fuller reference-data domain implementation.

\subsection{Runtime and transport}

Boost.Asio provides client and server transport over the shared envelope
protocol. The active server modes wire a local C++ \texttt{RaftNode} to an RPC
handler and transport server. The smoke CLI can also act as a client for Java or
C++ peers.

Use C++ smoke commands to validate boundary behavior:

\begin{lstlisting}[language=bash,caption={Selected C++ probe commands}]
graft-cpp/build/graft_smoke client-put 127.0.0.1 11082 k v1 cpp-node
graft-cpp/build/graft_smoke client-get 127.0.0.1 11082 k cpp-node
graft-cpp/build/graft_smoke cluster-summary 127.0.0.1 11082 cpp-cli
graft-cpp/build/graft_smoke reconfiguration-status 127.0.0.1 11082 cpp-cli
\end{lstlisting}

\subsection{Telemetry and observability}

The C++ implementation supports two layers of telemetry.

\subsubsection{In-cluster telemetry (protobuf RPC)}

All three telmetry RPC types are handled through the shared wire protocol:
\texttt{TelemetryRequest}/\texttt{TelemetryResponse},
\texttt{ClusterSummaryRequest}/\texttt{ClusterSummaryResponse}, and
\texttt{ReconfigurationStatusRequest}/\texttt{ReconfigurationStatusResponse}.
The \texttt{InMemoryRpcHandler} and \texttt{PersistentRpcHandler} classes in
\filepath{src/runtime/rpc\_handler.cpp} implement all three with authentication,
rate limiting, leader redirects and cluster health summaries matching the Java
reference behavior.

\subsubsection{External telemetry endpoints}

Configured via the \texttt{RAFT\_TELEMETRY\_EXPORTER} environment variable
(matching the Java \texttt{raft.telemetry.exporter} property):

\begin{itemize}
\item \texttt{prometheus}: starts a separate HTTP listener on
  \texttt{RAFT\_TELEMETRY\_PROMETHEUS\_HOST}:\texttt{RAFT\_TELEMETRY\_PROMETHEUS\_PORT}
  (default \texttt{127.0.0.1:9108}) and serves Raft metrics in Prometheus
  exposition format at \texttt{RAFT\_TELEMETRY\_PROMETHEUS\_PATH} (default
  \texttt{/metrics}). The listener is implemented directly on top of Boost.Asio
  without an external HTTP library.

\item \texttt{otlp} / \texttt{opentelemetry}: configuration is recognized and
  parsed but push export is not yet implemented (the timer skeleton is in
  place).

\item \texttt{otel-log} / \texttt{opentelemetry-log}: periodically writes a
  single-line JSON snapshot to \texttt{std::cout} at the configured
  \texttt{RAFT\_TELEMETRY\_EXPORT\_INTERVAL\_SECONDS} interval (default 15\,s).
  The JSON structure matches Java's \texttt{OpenTelemetryJsonExporter}.
\end{itemize}

The telemetry infrastructure lives in
\filepath{include/graft/telemetry/telemetry\_exporter.hpp} and
\filepath{src/telemetry/telemetry\_exporter.cpp}. The exporter is wired into
the active server modes in \filepath{src/app/cli.cpp}.

\subsection{Protobuf toolchain}

The C++ build uses \texttt{protoc} and \texttt{libprotobuf}. They should come
from the same package manager prefix. A warning such as compiler version
\texttt{33.x} versus library version \texttt{6.33.x} is usually a CMake
version-comparison artifact, but mixed prefixes can cause real compile or link
problems. Pin the compiler explicitly when needed:

\begin{lstlisting}[language=bash,caption={Pinning protoc}]
cmake -S graft-cpp -B graft-cpp/build \
  -DProtobuf_PROTOC_EXECUTABLE=/opt/local/bin/protoc
\end{lstlisting}

\subsection{Parity expectations}

C++ should be judged against Java at the Raft behavior boundary:

\begin{itemize}
\item status values such as \texttt{ACCEPTED}, \texttt{REDIRECT}, \texttt{RETRY}
and \texttt{INVALID},
\item vote and append semantics,
\item snapshot install and restore semantics,
\item read barrier and read-lease behavior,
\item joint-consensus majority rules,
\item membership admission and finalization,
\item mixed Java/C++ cluster behavior.
\end{itemize}

Differences in domain examples can be tolerated temporarily. Differences in
Raft safety semantics should be treated as defects.
